\subsection{An actual increase of an omitted-prime Frobenius order}

\begin{claim}{A.60}{Proved}{The cyclic-nineteen cover introduces an arithmetic factor of seven}
Let $C/\mathbb F _8$ and its pointed constant cyclic-nineteen
cover $D\to C$ be those of A.50, based at $P=(0,1)$.
Thus $C$ is the smooth proper model of
$y^2+x(x+1)y=x^5+1$, and $D$ is the pullback of the
degree-nineteen Lang quotient of its Jacobian along
$R\mapsto[R-P]$.
Take $\Sigma=\{3,19\}$ and use $\mathbb F _3$ coefficients
in the all-cover order test.
Then the prime-to-three parts of the actual Frobenius
orders on $H^1_{\mathrm{et}}$ of $C$ and $D$ are respectively
$4$ and $56$.  In particular,
\[
\begin{array}{c|cc}
 & C & D\\\hline
 v_7(\operatorname{ord}\operatorname{Frob}_8)&0&1\\
 v_2(\operatorname{ord}\operatorname{Frob}_8)&2&3.
\end{array}                                                   \tag{A.60.1}
\]
These are orders of the actual arithmetic action on a
geometric $\Sigma$-cover of one fixed curve over one fixed
finite field.
Consequently the all-cover supremum cannot in general be
replaced by the order on the base Jacobian alone.

For the canonical tower of A.59, alternating first the
full mod-$3$ and then the full mod-$19$ quotient, its second
level dominates this $D$.  Hence
\[
                         b_0(7)=0,\qquad b_2(7)\ge1.          \tag{A.60.2}
\]
Moreover, every further pointed geometrically connected
finite \'etale cover $E\to D$ over $\mathbb F _8$ with
geometric $3$-group Galois closure has
\[
 v_7\!\left(\operatorname{ord}
   (\operatorname{Frob}_8\mid H^1_{\mathrm{et}}
                     (E_{\overline{\mathbb F}_8},\mathbb F _3))\right)=1.
                                                                  \tag{A.60.3}
\]
Thus continuing with only $3$-group covers cannot turn
this factor $7$ into a factor $7^2$.
No unboundedness of the full alternating tower and no
classification in A12 is asserted.
For this curve and these primes, the explicit bound in
A.56 on the entire class of normal $3$-group by abelian
$19$-group covers is exactly one at the omitted prime
seven; this cover attains that bound.
\end{claim}

\begin{proof}
\medskip\noindent\textbf{The actual cover and its character polynomials.}
A.50 proves that the cover is geometrically connected,
finite \'etale, constant cyclic of degree nineteen, and
pointed over $\mathbb F _8$.  Its genus is twenty.
That proof also supplies the exact nontrivial character
factors.  If $\zeta$ is a primitive nineteenth root of
unity, put
\[
 S=\zeta+\zeta^{18}+\zeta^7+\zeta^{12}
                           +\zeta^8+\zeta^{11}.
\]
The factor belonging to one nontrivial character is
\[
                       L_\chi(T)=1+(3+S)T+8T^2.             \tag{A.60.4}
\]
The full list of nontrivial character factors replaces
$\zeta$ by its powers $\zeta^a$, $1\le a\le18$; there
are eighteen factors in total, including the displayed one.
Their product is the relative zeta numerator.
The sign in (A.60.4) is the sign of the Euler-product
coefficient in A.50.  Thus its point-count Frobenius
polynomial is $X^2+(3+S)X+8$.

The six exponents form the subgroup
$H=\langle8\rangle$ of $(\mathbb Z/19)^\times$.
There are three Gaussian periods, associated with its
three cosets, and each occurs six times in the eighteen
characters.  Reduction using
$1+\zeta+\cdots+\zeta^{18}=0$ gives
\[
                             S^3+S^2-6S-7=0.                \tag{A.60.5}
\]
This is also a finite exact identity in
$\mathbb Z[Z]/(1+Z+\cdots+Z^{18})$.
The cubic $f(U)=U^3+U^2-6U-7$ has no rational root:
the only possible integral roots are $\pm1,\pm7$, and
none is a root.  Hence it is irreducible, so those three
periods are precisely its three conjugates.

Their product of reciprocal character factors is
\[
\begin{split}
 N(X)&=\operatorname{Res}_U
             \bigl(f(U),X^2+(3+U)X+8\bigr)\\
     &=X^6+8X^5+39X^4+135X^3+312X^2+512X+512.
\end{split}                                                   \tag{A.60.6}
\]
For a direct expansion, write $A=X^2+3X+8$; then
$N=A^3-XA^2-6X^2A+7X^3$.
The repeated occurrence of the periods therefore gives
the exact primitive polynomial $N(X)^6$.

A.50 counts $9$ and $87$ points on $C$ over
$\mathbb F _8$ and $\mathbb F _{64}$, respectively.
The first trace is zero and the second is $-22$.
Newton's identities and the genus-two reciprocal
symmetry give
\[
 P_C(X)=X^4+11X^2+64,\qquad
                     P_D(X)=P_C(X)N(X)^6.                    \tag{A.60.7}
\]
As a consistency check, $N(1)=1519$ and
$1519^6=12284193340665686881$, the independently proved
relative class number in A.50.

\medskip\noindent\textbf{The exact modular orders.}
Over $\mathbb F _3$,
\[
 \overline P_C=(X^2+1)^2,
 \qquad n(X):=\overline N=X^6-X^5-X-1.                       \tag{A.60.8}
\]
The roots of $X^2+1$ have exact order four.
We give explicit finite identities for $n$, rather than
inferring its root orders from its degree.
In the quotient by $n$, the class $x$ of $X$ satisfies
\[
\begin{split}
 x^8&=x^5+x^3-x^2-x+1,\\
 x^{28}&=-1,\qquad x^{56}=1.
\end{split}                                                   \tag{A.60.9}
\]
In particular $x^{729}=x$, because $729=13\cdot56+1$.
The following two remainders have the indicated inverses
modulo $n$:
\[
\begin{array}{c|c}
 X^9-X\pmod n & X^5-X^4+X^2-X-1\\
 X^{27}-X\pmod n & -X^5-X^3+X+1
\end{array}                                                   \tag{A.60.10}
\]
Here the remainders in the left column are, explicitly,
\[
 X^5+X^4-X^3-X^2+X+1,
              \qquad -X^5+X^4-X+1,
\]
and multiplying each by its right-column entry gives
remainder one modulo $n$.
Thus $n$ divides $X^{3^6}-X$ and is coprime to both
$X^{3^2}-X$ and $X^{3^3}-X$.
Every irreducible factor has degree dividing six, and
the latter two coprimality statements exclude degrees
one, two and three.  Since $n$ itself has degree six,
it is irreducible.

Consequently its quotient is a field.  Equations
(A.60.9) show that the order of $x$ divides $56$, but
does not divide $28$ or $8$.  These are the quotients
of $56$ by its two prime divisors.  Hence its order is
exactly $56=2^3\cdot7$.

For a finite invertible matrix over $\mathbb F _3$,
the unipotent part has $3$-power order, whereas the
semisimple order is the least common multiple of the
orders of its eigenvalues.  Therefore (A.60.7)--(A.60.9)
give prime-to-three orders four and
$\operatorname{lcm}(4,56)=56$.
Arithmetic and point-count Frobenius are inverse
conventions on cohomology, as explained in A.59;
inversion preserves these finite orders.
This proves (A.60.1) for the actual arithmetic action,
independently of any possible Jordan blocks.

\medskip\noindent\textbf{The canonical tower and the next obstruction.}
Let $\Delta$ be the geometric pro-$\{3,19\}$ group and
$a:\Delta\twoheadrightarrow C_{19}$ the quotient of
the pointed cover $D$.
The image of $N_1=\Psi_3(\Delta)$ is all of $C_{19}$,
because cubing is an automorphism of that cyclic group.
The image of $N_2=\Psi_{19}(N_1)$ is trivial.
Thus $N_2\subseteq\ker a$, and the second pointed
canonical level maps to $D$ over the same field.
The rational-cohomology trace injection and the pro-$3$
congruence argument of A.59 give (A.60.2).

For completeness, put
$\Lambda=\pi_1(D_{\overline{\mathbb F}_8})^{(3)}$.
It is a finitely generated pro-$3$ group, and
\[
 \ker\bigl(\operatorname{Aut}(\Lambda)\longrightarrow
              \operatorname{Aut}(\Lambda/\Phi(\Lambda))\bigr)
                        \quad\hbox{is pro-}3.                \tag{A.60.11}
\]
Indeed take the characteristic open subgroups obtained
by intersecting the finitely many kernels of quotients
of each bounded order.  On each finite $3$-group quotient,
automorphisms trivial on its Frattini quotient form a
$3$-group: they act freely on the ordered generating
lifts of one generating tuple, whose number is a power
of three.  Triviality on the Frattini quotient of
$\Lambda$ implies that property on every finite quotient.
Passing to the inverse limit proves (A.60.11).

The image of the canonical $\mathbb Z_7$ factor of
$G_{\mathbb F _8}$ consequently has order exactly seven
on $\Lambda$, since it has that order on its mod-$3$
abelianization by (A.60.1), and meets the kernel in
(A.60.11) trivially.
Each pointed $3$-group cover of $D$ has its full induced
pro-$3$ topology, by the same finite-core argument as
in A.59.  Its based arithmetic action is the restriction
of the one on $\Lambda$, so its mod-$3$ cohomology has
$7$-part of order at most seven.
Conversely rational $3$-adic pullback from $D$ is injective
by trace, and reduction does not change the finite
$7$-part.  Its order is therefore at least seven.
This proves (A.60.3).

The stronger restricted-tower bound follows directly
from A.56.  Here $g=2$,
$\operatorname{ord}_{19}(3)=18$ and
$v_{19}(3^{18}-1)=1$, so its constant $c_{3,19}$ is one
and $M=\max\{4,2\}=4$.
Also $e_7(3)=v_7(3^6-1)=1$ and $v_7(19-1)=0$.
Finally $b_{19}(7)=0$: the group
$\operatorname{GL}_4(\mathbb F _{19})$ has order prime
to seven, since $\operatorname{ord}_7(19)=6>4$.
Thus A.56 gives
\[
                    B=0+1+0+\lfloor\log_7 4\rfloor=1.       \tag{A.60.12}
\]
The actual cover $D$ has trivial normal $3$-subgroup
and abelian $19$-quotient, so it belongs to that class
and (A.60.1) proves sharpness.

The finite calculation has been independently reproduced
in \texttt{a50\_cover19\_order\_jump.py},
with its retained JSON and output.  That implementation
checks the cyclotomic relation and conjugate multiplicities,
computes the norm both as a resultant and as a multiplication
determinant, and checks irreducibility and the exact modular
order.  Its geometric premise is precisely A.50.
The calculation is finite and does not test an infinite
family of covers.
\end{proof}

\paragraph{What an iteration would still have to prove.}
An orbit of length seven on an isotypic decomposition is
not by itself a factor of $7^2$ in Frobenius order.
If an invertible operator cyclically permutes $m$ nonzero
summands and its $m$th power restricts to $S$ on one
summand, then its order is $m\operatorname{ord}(S)$:
an identity power must first be divisible by $m$, and
the remaining condition is exactly an identity power
of $S$; the other restrictions are conjugate.
For $m=7$, a second factor of seven therefore requires
a factor of seven in this return action.
A.60 does not compute that return action on a further
alternating retained-prime cover.  Equation (A.60.3)
rules out obtaining the second factor by a further
pure pro-$3$ tower alone, and (A.60.12) requires any
successful construction to leave the entire class with
normal $3$-subgroup and abelian $19$-quotient.
